\section{南宋：行在的移动、军政的接续和江淮边界}

\begin{textwindow}{行在不是一个抽象的“南渡”}
《建炎以来系年要录》自建炎元年编排诏奏、行幸和军政消息；《宋史》保存相邻的本纪线索。它们最清楚地记录中枢如何移动，最不清楚的则是各地军民如何承受征发、逃亡和战乱。
\end{textwindow}

\subsection{一一二七--一一三二年：把朝廷重新接到道路和仓储上}
靖康之后，赵构即位。\eventanchor{SS-1127-EVENT-111}{高宗即位与南宋重建}这不是把旧朝廷完整搬到南方：扬州、镇江、杭州等地的驻跸选择，取决于河道、粮仓、勤王军和能够征调的税源。苗傅、刘正彦在杭州发动兵变，说明行在的禁军与宫廷关系尚未稳定；金军渡江及其北撤，则使江南水网和渡口成为双方行动的约束。

高宗朝廷后来以临安为较稳定的基地，并调整官署、度量衡和军政职能。这些政令能显示中枢试图把命令伸向各路，不能自动证明温台、江淮或四川以同一方式执行。南宋的形成因此是一连串接通粮道、军队和文书的处置，而非一场结束于一一二七年的迁徙。

\begin{turningpoint}{从逃离开封到经营临安}
临安的稳定让行在能够把税收、海运和官僚程序集中到东南；代价是北方故土、江淮与四川前线需要依赖更长的补给链。此后“守江”不只是军事判断，也是财政与交通的判断。
\end{turningpoint}

\subsection{绍兴和议与边界财政}
宋金反复作战、议和与再议和，最后形成以淮河--大散关一线为重的边界安排。和议为中枢腾出整理财政和军队的空间，却没有抹去淮西、两淮和四川的驻防压力。宋方诏奏所记录的盟约称谓，不能单独推出金方的目标；同样，金方材料中的战功也不能代替南宋地方社会的经验。

\begin{causalchain}{边界稳定为何仍需动员}
江南税粮和水路运输使朝廷能够维持行在与军队；前线的屯驻、转运和赏给反过来持续消耗这些资源。和议降低了大规模战事的频率，但未消除沿淮城镇、渡口和军户的成本。制度文本说明应当怎样征发，只有地方文书、碑刻和区域研究才可能检验局部落实。
\end{causalchain}

\subsection{海域、城市与晚期战局}
临安和泉州等地的市场、港口、造船及寺院网络，为南宋社会保留了多种收入和流动渠道；它们不应被简化成“富庶所以可守”。蒙古--元军向南推进时，襄阳、长江和江海航线又把这些资源转化为战争中的难题。南宋覆亡的不同地点和不同人群，须由宋、元、地方及海域材料分别追踪。

\begin{sourcenote}{可证实的序列与不可复原的经验}
本节以《宋史》、一一二七至一一六二年《建炎以来系年要录》及《宋会要辑稿》为年序和制度入口，并与金、元材料、契约、碑志和港口考古对读。它们不能给出全国人口流动、军队规模或社会反应的单一数字；这些空白保留为材料限度，而不以南宋中枢的声音补足。
\end{sourcenote}

\subsection{南宋年序：行在、江淮与海路的分段记录}

\paragraph{1127年（宋靖康二年；公元换算据《宋史》本纪纪年）：南宋}
\noindent \eventanchor{SYD-1127-SS-001}{癸未，觀文殿大學士唐恪仰藥自殺}\par

\paragraph{1127年（宋靖康二年；公元换算据《宋史》本纪纪年）：南宋}
\noindent \eventanchor{SYD-1127-SS-153}{何㮚、李若水勸帝親至軍中，從之，以太子監國而行}\par

\paragraph{1128年：南宋}
\noindent \eventanchor{SS-1128-EVENT-112}{高宗朝廷驻跸扬州等地，组织勤王与财政}\par

\paragraph{1128年（宋建炎二年；公元换算据《宋史》本纪纪年）：南宋}
\noindent \eventanchor{SYD-1128-SS-002}{罷殿中侍御史馬伸，尋責濮州}\par

\paragraph{1128年（宋建炎二年；公元换算据《宋史》本纪纪年）：南宋}
\noindent \eventanchor{SYD-1128-SS-154}{濱州賊蓋進陷棣州，守臣薑剛之死之}\par

\paragraph{1129年：南宋}
\noindent \eventanchor{SS-1129-EVENT-113}{苗傅、刘正彦发动兵变，随后被平定}\par

\paragraph{1129年（宋建炎三年；公元换算据《宋史》本纪纪年）：南宋}
\noindent \eventanchor{SYD-1129-SS-003}{丙午，粘罕陷徐州，守臣王復及子倚死之，軍校趙立結鄉兵為興復計}\par

\paragraph{1129年（宋建炎三年；公元换算据《宋史》本纪纪年）：南宋}
\noindent \eventanchor{SYD-1129-SS-155}{丙子，詔士民直言時政得失}\par

\paragraph{1130年：南宋与金}
\noindent \eventanchor{SS-1130-EVENT-114}{金军渡江后北撤，南宋保住江南核心地区}\par

\paragraph{1130年（宋建炎四年；公元换算据《宋史》本纪纪年）：南宋}
\noindent \eventanchor{SYD-1130-SS-004}{罷諸路帥臣兼制置使、諸州守臣兼管內安撫使}\par

\paragraph{1130年（宋建炎四年；公元换算据《宋史》本纪纪年）：南宋}
\noindent \eventanchor{SYD-1130-SS-156}{丙辰，復增左右司郎官為四員}\par

\paragraph{1131 年（共享年序桥接）：南宋、金与西夏（年序桥接）}
\noindent \eventanchor{SY-1131-CHRONOLOGY-BRIDGE}{【C级年序桥接】保留1131 年的多政权并行检索入口；本条不主张所列政权在当年发生直接互动，具体事件须另以单项材料入账。}\par

\paragraph{1131年（宋紹興元年；公元换算据《宋史》本纪纪年）：南宋}
\noindent \eventanchor{SYD-1131-SS-005}{丙辰，程昌寓遣杜湛擊楊華，敗之}\par

\paragraph{1131年（宋紹興元年；公元换算据《宋史》本纪纪年）：南宋}
\noindent \eventanchor{SYD-1131-SS-157}{丙申，遣內侍撫問孔彥舟、桑仲}\par

\paragraph{1132年：南宋}
\noindent \eventanchor{SS-1132-EVENT-115}{行在逐步以临安为稳定基地}\par

\paragraph{1132年（宋紹興二年；公元换算据《宋史》本纪纪年）：南宋}
\noindent \eventanchor{SYD-1132-SS-006}{班度量權衡于諸路，禁私造者}\par

\paragraph{1132年（宋紹興二年；公元换算据《宋史》本纪纪年）：南宋}
\noindent \eventanchor{SYD-1132-SS-158}{丙辰，禁溫、台二州民結集社會}\par

\paragraph{1133 年（共享年序桥接）：南宋、金与西夏（年序桥接）}
\noindent \eventanchor{SY-1133-CHRONOLOGY-BRIDGE}{【C级年序桥接】保留1133 年的多政权并行检索入口；本条不主张所列政权在当年发生直接互动，具体事件须另以单项材料入账。}\par

\paragraph{1133年（宋紹興三年；公元换算据《宋史》本纪纪年）：南宋}
\noindent \eventanchor{SYD-1133-SS-007}{安化蠻犯邊，廣西經略使許中發兵擊之}\par

\paragraph{1133年（宋紹興三年；公元换算据《宋史》本纪纪年）：南宋}
\noindent \eventanchor{SYD-1133-SS-159}{八月己丑，詔岳飛赴行在，留精兵萬人戍江州}\par

\paragraph{1134年：南宋与金}
\noindent \eventanchor{SS-1134-EVENT-116}{南宋军在襄汉、淮西等方向取得局部胜利}\par

\paragraph{1134年（宋紹興四年；公元换算据《宋史》本纪纪年）：南宋}
\noindent \eventanchor{SYD-1134-SS-008}{八月庚辰，以趙鼎知樞密院事，充川、陝宣撫處置使}\par

\paragraph{1134年（宋紹興四年；公元换算据《宋史》本纪纪年）：南宋}
\noindent \eventanchor{SYD-1134-SS-160}{編次建炎以來詔旨，頒諸路}\par

\paragraph{1135年：南宋}
\noindent \eventanchor{SS-1135-EVENT-117}{岳飞等将领获得更大经略与军队整编责任}\par

\paragraph{1135年（宋紹興五年；公元换算据《宋史》本纪纪年）：南宋}
1135年（宋紹興五年；公元换算据《宋史》本纪纪年） 的 《宋史》卷028〈紹興五年〉 记下 \eventanchor{SYD-1135-SS-009}{-\{岳\}-飛急攻湖賊水砦，賊將陳瑫降，楊太赴水死，餘黨劉衡等皆降}。对 南宋 而言，题名所点出的对象、机构或行动是这条记录能够落地的线索；本纪保留灾害或工程的报告地点与朝廷处置；灾情范围需另证。。它在后续年序中所打开或加重的是：可与地方文书、河工和环境材料比较，不将单点记录外推为全域因果。。《宋史》为元末官修，本纪保存宋廷纪年与处置；战果、数字、地方执行及对方动机须以编年、会要、对方史料或地方材料互校。

\paragraph{1135年（宋紹興五年；公元换算据《宋史》本纪纪年）：南宋}
\eventanchor{SYD-1135-SS-161}{-\{岳\}-飛為荊湖南北、襄陽府路制置使，將兵平湖賊楊太} 把 南宋 的行动固定在 1135年（宋紹興五年；公元换算据《宋史》本纪纪年），而 《宋史》卷028〈紹興五年〉 是核对该顺序的入口。这里首先可见的是题名中的人、物、城路或文书事项；本纪年条记录政令或机构处置；实施背景须参照制度材料。。直接的后续压力写在账本中为：可在同年及后续政令中追踪；条文不单独证明地方执行。。至于规模、地点细部或各方意图，《宋史》为元末官修，本纪保存宋廷纪年与处置；战果、数字、地方执行及对方动机须以编年、会要、对方史料或地方材料互校。

\paragraph{1136 年（共享年序桥接）：南宋、金与西夏（年序桥接）}
\noindent \eventanchor{SY-1136-CHRONOLOGY-BRIDGE}{【C级年序桥接】保留1136 年的多政权并行检索入口；本条不主张所列政权在当年发生直接互动，具体事件须另以单项材料入账。}\par

\paragraph{1136年（宋紹興六年；公元换算据《宋史》本纪纪年）：南宋}
\eventanchor{SYD-1136-SS-010}{-\{岳\}-飛及偽齊李成、孔彥舟連戰至蔡州，克之，偽守劉永壽舉城降} 所指的 1136年（宋紹興六年；公元换算据《宋史》本纪纪年） 节点，保留了 南宋 处理题名所涉人事、资源、地点或制度的痕迹；可复核的出处为 《宋史》卷028〈紹興六年〉。前一层条件不是宿命，而是 本纪从朝廷视角记录军政行动；出兵与战果需分辨。；随之可追踪的结果为 后续路线、控制与损失应与对方记录和遗址材料复核。。《宋史》为元末官修，本纪保存宋廷纪年与处置；战果、数字、地方执行及对方动机须以编年、会要、对方史料或地方材料互校。

\paragraph{1136年（宋紹興六年；公元换算据《宋史》本纪纪年）：南宋}
《宋史》卷028〈紹興六年〉 为 \eventanchor{SYD-1136-SS-162}{-\{岳\}-飛遣統制牛皋破偽齊鎮汝軍，禽其守薛亨} 提供年序凭据。南宋 在 1136年（宋紹興六年；公元换算据《宋史》本纪纪年） 面对的不是没有空间的政权名，而是题名所列的事务及其可调度的对象；本纪从朝廷视角记录军政行动；出兵与战果需分辨。。它使后续叙事必须继续追问 后续路线、控制与损失应与对方记录和遗址材料复核。。与此同时，《宋史》为元末官修，本纪保存宋廷纪年与处置；战果、数字、地方执行及对方动机须以编年、会要、对方史料或地方材料互校。

\paragraph{1137年（宋紹興七年；公元换算据《宋史》本纪纪年）：南宋}
1137年（宋紹興七年；公元换算据《宋史》本纪纪年），南宋 的 \eventanchor{SYD-1137-SS-011}{-\{岳\}-飛乞並統淮西兵以復京畿、陝右，許之，命飛盡護王德等諸將軍} 见于 《宋史》卷028〈紹興七年〉。题名留下的行动、文书、城路或人名构成此处唯一不应被抹平的材料细节；本纪年条记录政令或机构处置；实施背景须参照制度材料。。其影响只能谨慎写到 可在同年及后续政令中追踪；条文不单独证明地方执行。。《宋史》为元末官修，本纪保存宋廷纪年与处置；战果、数字、地方执行及对方动机须以编年、会要、对方史料或地方材料互校。

\paragraph{1137年（宋紹興七年；公元换算据《宋史》本纪纪年）：南宋}
\noindent \eventanchor{SYD-1137-SS-163}{罷淮南提點司，東西兩路各置轉運兼提點刑獄、提舉茶鹽常平事}\par

\paragraph{1138年：南宋}
\noindent \eventanchor{SS-1138-EVENT-118}{临安被确立为行在和事实都城}\par

\paragraph{1138年（宋紹興八年；公元换算据《宋史》本纪纪年）：南宋}
1138年（宋紹興八年；公元换算据《宋史》本纪纪年） 的 《宋史》卷029〈紹興八年〉 记下 \eventanchor{SYD-1138-SS-012}{從官張燾、晏敦復、魏矼、曾開、李彌遜、尹焞、梁汝嘉、樓炤、蘇符、薛徽言、-\{御\}-史方廷實皆言不可}。对 南宋 而言，题名所点出的对象、机构或行动是这条记录能够落地的线索；本纪年条记录政令或机构处置；实施背景须参照制度材料。。它在后续年序中所打开或加重的是：可在同年及后续政令中追踪；条文不单独证明地方执行。。《宋史》为元末官修，本纪保存宋廷纪年与处置；战果、数字、地方执行及对方动机须以编年、会要、对方史料或地方材料互校。

\paragraph{1138年（宋紹興八年；公元换算据《宋史》本纪纪年）：南宋}
\noindent \eventanchor{SYD-1138-SS-164}{八月戊午，詔：「日者遣使報聘鄰國，期還梓宮}\par

\paragraph{1139年（宋紹興九年；公元换算据《宋史》本纪纪年）：南宋}
\noindent \eventanchor{SYD-1139-SS-013}{八月己酉，復淮南諸州學官}\par

\paragraph{1139年（宋紹興九年；公元换算据《宋史》本纪纪年）：南宋}
\noindent \eventanchor{SYD-1139-SS-165}{丙申，命詳驗劉豫偽官，換給告身}\par

\paragraph{1140年：南宋与金}
\noindent \eventanchor{SS-1140-EVENT-119}{南宋军在河南、淮西等方向反攻}\par

\paragraph{1140年（宋紹興十年；公元换算据《宋史》本纪纪年）：南宋}
\noindent \eventanchor{SYD-1140-SS-014}{八月壬申朔，以張九成、喻樗、陳剛中、淩景夏、樊光遠、毛叔度、元盥等七人嘗不主和議，皆降黜之}\par

\paragraph{1140年（宋紹興十年；公元换算据《宋史》本纪纪年）：南宋}
\noindent \eventanchor{SYD-1140-SS-166}{丙午，命兩浙、江東、福建諸州團結弓弩手}\par

\paragraph{1141年（宋绍兴十一年、金皇统元年；干支、月日待核；公元换算据双方本纪年号对照）：南宋与金}
\noindent \eventanchor{SS-1141-EVENT-120}{南宋接受绍兴和议框架}\par

\paragraph{1141年（宋紹興十一年；公元换算据《宋史》本纪纪年）：南宋}
\noindent \eventanchor{SYD-1141-SS-015}{<u>李顯忠</u>遣統領<u>崔皋</u>擊敗<u>金</u>人于<u>舒城縣</u>}\par

\paragraph{1141年（宋紹興十一年；公元换算据《宋史》本纪纪年）：南宋}
\noindent \eventanchor{SYD-1141-SS-167}{<u>劉錡</u>自<u>東關</u>擊敗<u>金</u>人于<u>青谿</u>}\par

\paragraph{1142年：南宋}
\noindent \eventanchor{SS-1142-EVENT-121}{岳飞被处死，相关将领和军队受到整顿}\par

\paragraph{1142年（宋紹興十二年；公元换算据《宋史》本纪纪年）：南宋}
\noindent \eventanchor{SYD-1142-SS-016}{八月辛酉朔，兀朮使來求商州及和尚、方山二原}\par

\paragraph{1142年（宋紹興十二年；公元换算据《宋史》本纪纪年）：南宋}
\noindent \eventanchor{SYD-1142-SS-168}{丙午，金使劉筈、完顏宗表等九人入見}\par

\paragraph{1143 年（共享年序桥接）：南宋、金与西夏（年序桥接）}
\noindent \eventanchor{SY-1143-CHRONOLOGY-BRIDGE}{【C级年序桥接】保留1143 年的多政权并行检索入口；本条不主张所列政权在当年发生直接互动，具体事件须另以单项材料入账。}\par

\paragraph{1143年（宋紹興十三年；公元换算据《宋史》本纪纪年）：南宋}
\noindent \eventanchor{SYD-1143-SS-017}{八月丙戌，遣吏部侍郎江邈奉迎累朝神禦於溫州}\par

\paragraph{1143年（宋紹興十三年；公元换算据《宋史》本纪纪年）：南宋}
\noindent \eventanchor{SYD-1143-SS-169}{丙寅，處州兵士楊興等謀作亂，事覺伏誅}\par

\paragraph{1144年（宋紹興十四年；公元换算据《宋史》本纪纪年）：南宋}
1144年（宋紹興十四年；公元换算据《宋史》本纪纪年） 的 《宋史》卷030〈紹興十四年〉 记下 \eventanchor{SYD-1144-SS-018}{丙午，罷-\{万俟\}-禼}。对 南宋 而言，题名所点出的对象、机构或行动是这条记录能够落地的线索；本纪年条记录政令或机构处置；实施背景须参照制度材料。。它在后续年序中所打开或加重的是：可在同年及后续政令中追踪；条文不单独证明地方执行。。《宋史》为元末官修，本纪保存宋廷纪年与处置；战果、数字、地方执行及对方动机须以编年、会要、对方史料或地方材料互校。

\paragraph{1144年（宋紹興十四年；公元换算据《宋史》本纪纪年）：南宋}
\noindent \eventanchor{SYD-1144-SS-170}{丙寅，立明法科兼經法}\par

\paragraph{1145 年（共享年序桥接）：南宋、金与西夏（年序桥接）}
\noindent \eventanchor{SY-1145-CHRONOLOGY-BRIDGE}{【C级年序桥接】保留1145 年的多政权并行检索入口；本条不主张所列政权在当年发生直接互动，具体事件须另以单项材料入账。}\par

\paragraph{1145年（宋紹興十五年；公元换算据《宋史》本纪纪年）：南宋}
\noindent \eventanchor{SYD-1145-SS-019}{八月申戌朔，禁收折帛合零錢，止輸實數}\par

\paragraph{1145年（宋紹興十五年；公元换算据《宋史》本纪纪年）：南宋}
\noindent \eventanchor{SYD-1145-SS-171}{丙寅，全給秦檜歲賜公使錢萬緡}\par

\paragraph{1146 年（共享年序桥接）：南宋、金与西夏（年序桥接）}
\noindent \eventanchor{SY-1146-CHRONOLOGY-BRIDGE}{【C级年序桥接】保留1146 年的多政权并行检索入口；本条不主张所列政权在当年发生直接互动，具体事件须另以单项材料入账。}\par

\paragraph{1146年（宋紹興十六年；公元换算据《宋史》本纪纪年）：南宋}
\noindent \eventanchor{SYD-1146-SS-020}{丙申，復何鑄為端明殿學士兼侍讀}\par

\paragraph{1146年（宋紹興十六年；公元换算据《宋史》本纪纪年）：南宋}
\noindent \eventanchor{SYD-1146-SS-172}{丁亥，金遣烏古論海等來賀天申節}\par

\paragraph{1147 年（共享年序桥接）：南宋、金与西夏（年序桥接）}
\noindent \eventanchor{SY-1147-CHRONOLOGY-BRIDGE}{【C级年序桥接】保留1147 年的多政权并行检索入口；本条不主张所列政权在当年发生直接互动，具体事件须另以单项材料入账。}\par

\paragraph{1147年（宋紹興十七年；公元换算据《宋史》本纪纪年）：南宋}
\noindent \eventanchor{SYD-1147-SS-021}{八月庚子，罷建州創置賣鹽坊}\par

\paragraph{1147年（宋紹興十七年；公元换算据《宋史》本纪纪年）：南宋}
\noindent \eventanchor{SYD-1147-SS-173}{丙辰，金遣完顏宗藩等來賀明年正旦}\par

\paragraph{1148 年（共享年序桥接）：南宋、金与西夏（年序桥接）}
\noindent \eventanchor{SY-1148-CHRONOLOGY-BRIDGE}{【C级年序桥接】保留1148 年的多政权并行检索入口；本条不主张所列政权在当年发生直接互动，具体事件须另以单项材料入账。}\par

\paragraph{1148年（宋紹興十八年；公元换算据《宋史》本纪纪年）：南宋}
\noindent \eventanchor{SYD-1148-SS-022}{丙辰，加士夽開府儀同三司}\par

\paragraph{1148年（宋紹興十八年；公元换算据《宋史》本纪纪年）：南宋}
\noindent \eventanchor{SYD-1148-SS-174}{丙寅，借給被災農民春耕費}\par

\paragraph{1149年（宋紹興十九年；公元换算据《宋史》本纪纪年）：南宋}
\noindent \eventanchor{SYD-1149-SS-023}{丁丑，金遣完顏袞等來賀明年正旦}\par

\paragraph{1149年（宋紹興十九年；公元换算据《宋史》本纪纪年）：南宋}
\noindent \eventanchor{SYD-1149-SS-175}{丁未，減連、英、循、惠、新、恩六州免行錢}\par

\paragraph{1150年：南宋}
\noindent \eventanchor{SS-1150-EVENT-122}{朝廷调整会子与纸币管理}\par

\paragraph{1150年（宋紹興二十年；公元换算据《宋史》本纪纪年）：南宋}
\noindent \eventanchor{SYD-1150-SS-024}{丙申，侍御史曹筠以附下罔上罷}\par

\paragraph{1150年（宋紹興二十年；公元换算据《宋史》本纪纪年）：南宋}
\noindent \eventanchor{SYD-1150-SS-176}{丙午，兩浙轉運副使曹泳言，李孟堅誦其父光所撰私史，語涉譏謗，詔送大理寺}\par

\paragraph{1151 年（共享年序桥接）：南宋、金与西夏（年序桥接）}
\noindent \eventanchor{SY-1151-CHRONOLOGY-BRIDGE}{【C级年序桥接】保留1151 年的多政权并行检索入口；本条不主张所列政权在当年发生直接互动，具体事件须另以单项材料入账。}\par

\paragraph{1151年（宋紹興二十一年；公元换算据《宋史》本纪纪年）：南宋}
\noindent \eventanchor{SYD-1151-SS-025}{八月辛未，秦檜上《重修諸路茶鹽法》}\par

\paragraph{1151年（宋紹興二十一年；公元换算据《宋史》本纪纪年）：南宋}
\noindent \eventanchor{SYD-1151-SS-177}{丁亥，賜禮部進士趙逵以下四百四人及第、出身}\par

\paragraph{1152 年（共享年序桥接）：南宋、金与西夏（年序桥接）}
\noindent \eventanchor{SY-1152-CHRONOLOGY-BRIDGE}{【C级年序桥接】保留1152 年的多政权并行检索入口；本条不主张所列政权在当年发生直接互动，具体事件须另以单项材料入账。}\par

\paragraph{1152年（宋紹興二十二年；公元换算据《宋史》本纪纪年）：南宋}
\noindent \eventanchor{SYD-1152-SS-026}{八月己卯，遣鄂州都統制田師中發兵同江西安撫使張澄、殿前司游奕軍統制李耕討述}\par

\paragraph{1152年（宋紹興二十二年；公元换算据《宋史》本纪纪年）：南宋}
\noindent \eventanchor{SYD-1152-SS-178}{丁巳，立薦舉受財刑名}\par

\paragraph{1153年（宋紹興二十三年；公元换算据《宋史》本纪纪年）：南宋}
\noindent \eventanchor{SYD-1153-SS-027}{八月乙丑，士撙薨，追封韶王}\par

\paragraph{1153年（宋紹興二十三年；公元换算据《宋史》本纪纪年）：南宋}
\noindent \eventanchor{SYD-1153-SS-179}{丙寅，左宣教郎王孝廉謀據成都叛，事覺，伏誅}\par

\paragraph{1154 年（共享年序桥接）：南宋、金与西夏（年序桥接）}
\noindent \eventanchor{SY-1154-CHRONOLOGY-BRIDGE}{【C级年序桥接】保留1154 年的多政权并行检索入口；本条不主张所列政权在当年发生直接互动，具体事件须另以单项材料入账。}\par

\paragraph{1154年（宋紹興二十四年；公元换算据《宋史》本纪纪年）：南宋}
\noindent \eventanchor{SYD-1154-SS-028}{尋遣鐘世明如四川同議}\par

\paragraph{1154年（宋紹興二十四年；公元换算据《宋史》本纪纪年）：南宋}
\noindent \eventanchor{SYD-1154-SS-180}{八月壬辰，禁百官避免輪對}\par

\paragraph{1155 年（共享年序桥接）：南宋、金与西夏（年序桥接）}
\noindent \eventanchor{SY-1155-CHRONOLOGY-BRIDGE}{【C级年序桥接】保留1155 年的多政权并行检索入口；本条不主张所列政权在当年发生直接互动，具体事件须另以单项材料入账。}\par

\paragraph{1155年（宋紹興二十五年；公元换算据《宋史》本纪纪年）：南宋}
\eventanchor{SYD-1155-SS-029}{詔張浚、折彥質、-\{万俟\}-禼、段拂聽自便} 的可见载体是 《宋史》卷031〈紹興二十五年〉；它把 南宋 放在 1155年（宋紹興二十五年；公元换算据《宋史》本纪纪年） 的这一处置里。本纪年条记录政令或机构处置；实施背景须参照制度材料。 所以这不是可由相邻年份自动推出的结果；其后留下的条件是：可在同年及后续政令中追踪；条文不单独证明地方执行。。材料边界也须保留：《宋史》为元末官修，本纪保存宋廷纪年与处置；战果、数字、地方执行及对方动机须以编年、会要、对方史料或地方材料互校。

\paragraph{1155年（宋紹興二十五年；公元换算据《宋史》本纪纪年）：南宋}
\noindent \eventanchor{SYD-1155-SS-181}{丙申，復以蕭振為四川制置使}\par

\paragraph{1156年（宋紹興二十六年；公元换算据《宋史》本纪纪年）：南宋}
\noindent \eventanchor{SYD-1156-SS-030}{八月戊寅，班元豐、崇甯學制于諸路}\par

\paragraph{1156年（宋紹興二十六年；公元换算据《宋史》本纪纪年）：南宋}
\noindent \eventanchor{SYD-1156-SS-182}{丙午，立互易薦舉坐罪法}\par

\paragraph{1157 年（共享年序桥接）：南宋、金与西夏（年序桥接）}
\noindent \eventanchor{SY-1157-CHRONOLOGY-BRIDGE}{【C级年序桥接】保留1157 年的多政权并行检索入口；本条不主张所列政权在当年发生直接互动，具体事件须另以单项材料入账。}\par

\paragraph{1157年（宋紹興二十七年；公元换算据《宋史》本纪纪年）：南宋}
\noindent \eventanchor{SYD-1157-SS-031}{丙辰，初命州縣置禁曆}\par

\paragraph{1157年（宋紹興二十七年；公元换算据《宋史》本纪纪年）：南宋}
\noindent \eventanchor{SYD-1157-SS-183}{丙戌，賜禮部進士王十朋以下四百二十六人及第、出身}\par

\paragraph{1158 年（共享年序桥接）：南宋、金与西夏（年序桥接）}
\noindent \eventanchor{SY-1158-CHRONOLOGY-BRIDGE}{【C级年序桥接】保留1158 年的多政权并行检索入口；本条不主张所列政权在当年发生直接互动，具体事件须另以单项材料入账。}\par

\paragraph{1158年（宋紹興二十八年；公元换算据《宋史》本纪纪年）：南宋}
\noindent \eventanchor{SYD-1158-SS-032}{己卯，命取公私銅器悉付鑄錢司，民間不輸者罪之}\par

\paragraph{1158年（宋紹興二十八年；公元换算据《宋史》本纪纪年）：南宋}
\noindent \eventanchor{SYD-1158-SS-184}{八月戊子朔，置國史院，修神、哲、徽三朝正史}\par

\paragraph{1159 年（共享年序桥接）：南宋、金与西夏（年序桥接）}
\noindent \eventanchor{SY-1159-CHRONOLOGY-BRIDGE}{【C级年序桥接】保留1159 年的多政权并行检索入口；本条不主张所列政权在当年发生直接互动，具体事件须另以单项材料入账。}\par

\paragraph{1159年（宋紹興二十九年；公元换算据《宋史》本纪纪年）：南宋}
\noindent \eventanchor{SYD-1159-SS-033}{八月甲子，募商人輸米行在諸倉，願以茶、鹽、礬鈔等償直者聽}\par

\paragraph{1159年（宋紹興二十九年；公元换算据《宋史》本纪纪年）：南宋}
\noindent \eventanchor{SYD-1159-SS-185}{丙午，禁內外將佐營造、回易，掊斂軍士}\par

\paragraph{1160 年（共享年序桥接）：南宋、金与西夏（年序桥接）}
\noindent \eventanchor{SY-1160-CHRONOLOGY-BRIDGE}{【C级年序桥接】保留1160 年的多政权并行检索入口；本条不主张所列政权在当年发生直接互动，具体事件须另以单项材料入账。}\par

\paragraph{1160年（宋紹興三十年；公元换算据《宋史》本纪纪年）：南宋}
\noindent \eventanchor{SYD-1160-SS-034}{丙申，金遣蕭榮等來賀天申節}\par

\paragraph{1160年（宋紹興三十年；公元换算据《宋史》本纪纪年）：南宋}
\noindent \eventanchor{SYD-1160-SS-186}{丙申，命劉寶招制勝軍千人}\par

\paragraph{1161年（宋绍兴三十一年、金正隆六年；干支、月日待核；公元换算据双方本纪年号对照）：南宋与金}
\noindent \eventanchor{SS-1161-EVENT-123}{南宋在采石等地击退金军渡江攻势}\par

\paragraph{1161年（宋紹興三十一年；公元换算据《宋史》本纪纪年）：南宋}
\noindent \eventanchor{SYD-1161-SS-035}{乙未，行新造會子於淮、浙、湖北、京西諸州}\par

\paragraph{1161年（宋紹興三十一年；公元换算据《宋史》本纪纪年）：南宋}
\noindent \eventanchor{SYD-1161-SS-187}{八月辛丑朔，忠義人魏勝復海州，李寶承制以勝知州事}\par

\paragraph{1162年：南宋}
\noindent \eventanchor{SS-1162-EVENT-124}{高宗禅位，孝宗即位}\par

\paragraph{1162年（宋紹興三十二年；公元换算据《宋史》本纪纪年）：南宋}
\noindent \eventanchor{SYD-1162-SS-036}{丙戌，給張浚錢十九萬緡，造沿江諸軍戰艦}\par

\paragraph{1162年（宋紹興三十二年；公元换算据《宋史》本纪纪年）：南宋}
\noindent \eventanchor{SYD-1162-SS-188}{丙子，姚仲遣將復原州}\par

\paragraph{1163年：南宋与金}
\noindent \eventanchor{SS-1163-EVENT-125}{孝宗支持北伐，宋军在符离等地受挫}\par

\paragraph{1163年（宋隆興元年；公元换算据《宋史》本纪纪年）：南宋}
\noindent \eventanchor{SYD-1163-SS-037}{安慶軍節度使士篯乞減奉賜之半，以助軍用}\par

\paragraph{1163年（宋隆興元年；公元换算据《宋史》本纪纪年）：南宋}
\noindent \eventanchor{SYD-1163-SS-189}{八月丙寅，張浚復都督江、淮軍馬}\par

\paragraph{1164年（宋隆兴二年、金大定四年；干支、月日待核；公元换算据双方本纪年号对照）：南宋与金}
\noindent \eventanchor{SS-1164-EVENT-126}{宋金订立隆兴和议}\par

\paragraph{1164年（宋隆興二年；公元换算据《宋史》本纪纪年）：南宋}
\noindent \eventanchor{SYD-1164-SS-038}{罷兩浙、福建、江西、湖南、夔州路參議官}\par

\paragraph{1164年（宋隆興二年；公元换算据《宋史》本纪纪年）：南宋}
\noindent \eventanchor{SYD-1164-SS-190}{丙申，命虞允文調兵討廣西諸盜}\par

\paragraph{1165 年（共享年序桥接）：南宋、金与西夏（年序桥接）}
\noindent \eventanchor{SY-1165-CHRONOLOGY-BRIDGE}{【C级年序桥接】保留1165 年的多政权并行检索入口；本条不主张所列政权在当年发生直接互动，具体事件须另以单项材料入账。}\par

\paragraph{1165年（宋乾道元年；公元换算据《宋史》本纪纪年）：南宋}
\noindent \eventanchor{SYD-1165-SS-039}{八月己卯，以永豐圩田賜建康都統司}\par

\paragraph{1165年（宋乾道元年；公元换算据《宋史》本纪纪年）：南宋}
\noindent \eventanchor{SYD-1165-SS-191}{丙辰，詔有司治皇后家廟}\par

\paragraph{1166 年（共享年序桥接）：南宋、金与西夏（年序桥接）}
\noindent \eventanchor{SY-1166-CHRONOLOGY-BRIDGE}{【C级年序桥接】保留1166 年的多政权并行检索入口；本条不主张所列政权在当年发生直接互动，具体事件须另以单项材料入账。}\par

\paragraph{1166年（宋乾道二年；公元换算据《宋史》本纪纪年）：南宋}
\noindent \eventanchor{SYD-1166-SS-040}{丁卯，賜禮部進士<u>蕭國梁</u>以下四百九十有三人及第、出身}\par

\paragraph{1166年（宋乾道二年；公元换算据《宋史》本纪纪年）：南宋}
\noindent \eventanchor{SYD-1166-SS-192}{八月辛未朔，詔兩淮行鐵錢，銅錢毋過江北}\par

\paragraph{1167 年（共享年序桥接）：南宋、金与西夏（年序桥接）}
\noindent \eventanchor{SY-1167-CHRONOLOGY-BRIDGE}{【C级年序桥接】保留1167 年的多政权并行检索入口；本条不主张所列政权在当年发生直接互动，具体事件须另以单项材料入账。}\par

\paragraph{1167年（宋乾道三年；公元换算据《宋史》本纪纪年）：南宋}
\noindent \eventanchor{SYD-1167-SS-041}{八月丁酉，內侍陳瑜、李宗回坐交結戚方受賂，瑜除名、決杖、黥面配循州，宗回除名、筠州編管，方責授果州團練副使、潭州安置，籍所盜庫金以犒軍}\par

\paragraph{1167年（宋乾道三年；公元换算据《宋史》本纪纪年）：南宋}
\noindent \eventanchor{SYD-1167-SS-193}{罷成都、潼川路轉運司輪年銓試，以其事付制置司}\par

\paragraph{1168 年（共享年序桥接）：南宋、金与西夏（年序桥接）}
\noindent \eventanchor{SY-1168-CHRONOLOGY-BRIDGE}{【C级年序桥接】保留1168 年的多政权并行检索入口；本条不主张所列政权在当年发生直接互动，具体事件须另以单项材料入账。}\par

\paragraph{1168年（宋乾道四年；公元换算据《宋史》本纪纪年）：南宋}
\noindent \eventanchor{SYD-1168-SS-042}{八月乙未，班祈雨雪之法于諸路}\par

\paragraph{1168年（宋乾道四年；公元换算据《宋史》本纪纪年）：南宋}
\noindent \eventanchor{SYD-1168-SS-194}{丁亥，以饒、信二州、建甯府饑民嘯聚，遣官措置振濟}\par

\paragraph{1169 年（共享年序桥接）：南宋、金与西夏（年序桥接）}
\noindent \eventanchor{SY-1169-CHRONOLOGY-BRIDGE}{【C级年序桥接】保留1169 年的多政权并行检索入口；本条不主张所列政权在当年发生直接互动，具体事件须另以单项材料入账。}\par

\paragraph{1169年（宋乾道五年；公元换算据《宋史》本纪纪年）：南宋}
\noindent \eventanchor{SYD-1169-SS-043}{罷利州路諸州營田官兵，募民耕佃}\par

\paragraph{1169年（宋乾道五年；公元换算据《宋史》本纪纪年）：南宋}
\noindent \eventanchor{SYD-1169-SS-195}{丙寅，爲岳飛立廟于鄂州}\par

\paragraph{1170 年（共享年序桥接）：南宋、金与西夏（年序桥接）}
\noindent \eventanchor{SY-1170-CHRONOLOGY-BRIDGE}{【C级年序桥接】保留1170 年的多政权并行检索入口；本条不主张所列政权在当年发生直接互动，具体事件须另以单项材料入账。}\par

\paragraph{1170年（宋乾道六年；公元换算据《宋史》本纪纪年）：南宋}
\noindent \eventanchor{SYD-1170-SS-044}{八月庚戌，虞允文請蚤建太子}\par

\paragraph{1170年（宋乾道六年；公元换算据《宋史》本纪纪年）：南宋}
\noindent \eventanchor{SYD-1170-SS-196}{罷行在至鎮江徵稅所比近者十有三}\par

\paragraph{1171 年（共享年序桥接）：南宋、金与西夏（年序桥接）}
\noindent \eventanchor{SY-1171-CHRONOLOGY-BRIDGE}{【C级年序桥接】保留1171 年的多政权并行检索入口；本条不主张所列政权在当年发生直接互动，具体事件须另以单项材料入账。}\par

\paragraph{1171年（宋乾道七年；公元换算据《宋史》本纪纪年）：南宋}
\noindent \eventanchor{SYD-1171-SS-045}{八月丙辰，詔兩淮民丁充民兵者，本名丁錢勿輸}\par

\paragraph{1171年（宋乾道七年；公元换算据《宋史》本纪纪年）：南宋}
\noindent \eventanchor{SYD-1171-SS-197}{丙寅，金遣完顏宗甯等來賀明年正旦}\par

\paragraph{1172 年（共享年序桥接）：南宋、金与西夏（年序桥接）}
\noindent \eventanchor{SY-1172-CHRONOLOGY-BRIDGE}{【C级年序桥接】保留1172 年的多政权并行检索入口；本条不主张所列政权在当年发生直接互动，具体事件须另以单项材料入账。}\par

\paragraph{1172年（宋乾道八年；公元换算据《宋史》本纪纪年）：南宋}
\noindent \eventanchor{SYD-1172-SS-046}{八年春正月庚午朔，班《乾道敕令格式》}\par

\paragraph{1172年（宋乾道八年；公元换算据《宋史》本纪纪年）：南宋}
\noindent \eventanchor{SYD-1172-SS-198}{丙辰，金遣夾穀清臣等來賀會慶節}\par

\paragraph{1173年：南宋}
\noindent \eventanchor{SS-1173-EVENT-127}{朱熹等在书院、讲学与地方事务中活动}\par

\paragraph{1173年（宋乾道九年；公元换算据《宋史》本纪纪年）：南宋}
\noindent \eventanchor{SYD-1173-SS-047}{八月丙子，詔興修水利}\par

\paragraph{1173年（宋乾道九年；公元换算据《宋史》本纪纪年）：南宋}
\noindent \eventanchor{SYD-1173-SS-199}{丙辰，復分淮南安撫司爲東、西路}\par

\paragraph{1174 年（共享年序桥接）：南宋、金与西夏（年序桥接）}
\noindent \eventanchor{SY-1174-CHRONOLOGY-BRIDGE}{【C级年序桥接】保留1174 年的多政权并行检索入口；本条不主张所列政权在当年发生直接互动，具体事件须另以单项材料入账。}\par

\paragraph{1174年（宋淳熙元年；公元换算据《宋史》本纪纪年）：南宋}
\noindent \eventanchor{SYD-1174-SS-048}{八月己未，張說罷爲太尉，提舉隆興府玉隆觀}\par

\paragraph{1174年（宋淳熙元年；公元换算据《宋史》本纪纪年）：南宋}
\noindent \eventanchor{SYD-1174-SS-200}{丙午，禁兩淮耕牛出境}\par

\paragraph{1175 年（共享年序桥接）：南宋、金与西夏（年序桥接）}
\noindent \eventanchor{SY-1175-CHRONOLOGY-BRIDGE}{【C级年序桥接】保留1175 年的多政权并行检索入口；本条不主张所列政权在当年发生直接互动，具体事件须另以单项材料入账。}\par

\paragraph{1175年（宋淳熙二年；公元换算据《宋史》本纪纪年）：南宋}
\noindent \eventanchor{SYD-1175-SS-049}{八月丙辰，江西總管賈和仲以捕茶寇失律，除名、賀州編管}\par

\paragraph{1175年（宋淳熙二年；公元换算据《宋史》本纪纪年）：南宋}
\eventanchor{SYD-1175-SS-201}{丁-\{丑\}-，遣左司諫湯邦彥等使金申議} 所指的 1175年（宋淳熙二年；公元换算据《宋史》本纪纪年） 节点，保留了 南宋 处理题名所涉人事、资源、地点或制度的痕迹；可复核的出处为 《宋史》卷034〈淳熙二年〉。前一层条件不是宿命，而是 本纪记录使节、贡赐或往来，不等于稳定统属关系。；随之可追踪的结果为 可与对方编年和交通、文书材料核对其后续落实。。《宋史》为元末官修，本纪保存宋廷纪年与处置；战果、数字、地方执行及对方动机须以编年、会要、对方史料或地方材料互校。

\paragraph{1176 年（共享年序桥接）：南宋、金与西夏（年序桥接）}
\noindent \eventanchor{SY-1176-CHRONOLOGY-BRIDGE}{【C级年序桥接】保留1176 年的多政权并行检索入口；本条不主张所列政权在当年发生直接互动，具体事件须另以单项材料入账。}\par

\paragraph{1176年（宋淳熙三年；公元换算据《宋史》本纪纪年）：南宋}
1176年（宋淳熙三年；公元换算据《宋史》本纪纪年），南宋 的 \eventanchor{SYD-1176-SS-050}{丁-\{丑\}-，命臨安守臣嚴禁逾侈} 见于 《宋史》卷034〈淳熙三年〉。题名留下的行动、文书、城路或人名构成此处唯一不应被抹平的材料细节；本纪年条记录政令或机构处置；实施背景须参照制度材料。。其影响只能谨慎写到 可在同年及后续政令中追踪；条文不单独证明地方执行。。《宋史》为元末官修，本纪保存宋廷纪年与处置；战果、数字、地方执行及对方动机须以编年、会要、对方史料或地方材料互校。

\paragraph{1176年（宋淳熙三年；公元换算据《宋史》本纪纪年）：南宋}
\noindent \eventanchor{SYD-1176-SS-202}{丁酉，湯邦彥、陳雷奉使無狀，除名，邦彥新州、雷永州編管}\par

\paragraph{1177 年（共享年序桥接）：南宋、金与西夏（年序桥接）}
\noindent \eventanchor{SY-1177-CHRONOLOGY-BRIDGE}{【C级年序桥接】保留1177 年的多政权并行检索入口；本条不主张所列政权在当年发生直接互动，具体事件须另以单项材料入账。}\par

\paragraph{1177年（宋淳熙四年；公元换算据《宋史》本纪纪年）：南宋}
\noindent \eventanchor{SYD-1177-SS-051}{八月辛巳，禁耕牛過淮}\par

\paragraph{1177年（宋淳熙四年；公元换算据《宋史》本纪纪年）：南宋}
\eventanchor{SYD-1177-SS-203}{丁-\{丑\}-，詔監司、守臣歲舉武臣堪知縣者各二人} 把 南宋 的行动固定在 1177年（宋淳熙四年；公元换算据《宋史》本纪纪年），而 《宋史》卷034〈淳熙四年〉 是核对该顺序的入口。这里首先可见的是题名中的人、物、城路或文书事项；本纪年条记录政令或机构处置；实施背景须参照制度材料。。直接的后续压力写在账本中为：可在同年及后续政令中追踪；条文不单独证明地方执行。。至于规模、地点细部或各方意图，《宋史》为元末官修，本纪保存宋廷纪年与处置；战果、数字、地方执行及对方动机须以编年、会要、对方史料或地方材料互校。

\paragraph{1178 年（共享年序桥接）：南宋、金与西夏（年序桥接）}
\noindent \eventanchor{SY-1178-CHRONOLOGY-BRIDGE}{【C级年序桥接】保留1178 年的多政权并行检索入口；本条不主张所列政权在当年发生直接互动，具体事件须另以单项材料入账。}\par

\paragraph{1178年（宋淳熙五年；公元换算据《宋史》本纪纪年）：南宋}
\noindent \eventanchor{SYD-1178-SS-052}{八月甲午，詔諸路監司戒所部，民稅毋以重價強折輸錢}\par

\paragraph{1178年（宋淳熙五年；公元换算据《宋史》本纪纪年）：南宋}
\noindent \eventanchor{SYD-1178-SS-204}{丙辰，金遣烏延察等來賀明年正旦}\par

\paragraph{1179 年（共享年序桥接）：南宋、金与西夏（年序桥接）}
\noindent \eventanchor{SY-1179-CHRONOLOGY-BRIDGE}{【C级年序桥接】保留1179 年的多政权并行检索入口；本条不主张所列政权在当年发生直接互动，具体事件须另以单项材料入账。}\par

\paragraph{1179年（宋淳熙六年；公元换算据《宋史》本纪纪年）：南宋}
\noindent \eventanchor{SYD-1179-SS-053}{八月庚寅，罷諸路監司、帥守便宜行事}\par

\paragraph{1179年（宋淳熙六年；公元换算据《宋史》本纪纪年）：南宋}
\noindent \eventanchor{SYD-1179-SS-205}{丙申，詔太學兩優釋褐，與殿試第二人恩例}\par

\paragraph{1180 年（共享年序桥接）：南宋、金与西夏（年序桥接）}
\noindent \eventanchor{SY-1180-CHRONOLOGY-BRIDGE}{【C级年序桥接】保留1180 年的多政权并行检索入口；本条不主张所列政权在当年发生直接互动，具体事件须另以单项材料入账。}\par

\paragraph{1180年（宋淳熙七年；公元换算据《宋史》本纪纪年）：南宋}
\noindent \eventanchor{SYD-1180-SS-054}{戊子，遣葉宏等使金賀正旦}\par

\paragraph{1180年（宋淳熙七年；公元换算据《宋史》本纪纪年）：南宋}
\noindent \eventanchor{SYD-1180-SS-206}{八月癸未，禁黎州官吏市蕃商物}\par

\paragraph{1181 年（共享年序桥接）：南宋、金与西夏（年序桥接）}
\noindent \eventanchor{SY-1181-CHRONOLOGY-BRIDGE}{【C级年序桥接】保留1181 年的多政权并行检索入口；本条不主张所列政权在当年发生直接互动，具体事件须另以单项材料入账。}\par

\paragraph{1181年（宋淳熙八年；公元换算据《宋史》本纪纪年）：南宋}
\noindent \eventanchor{SYD-1181-SS-055}{八月丙午，以旱，罷招軍}\par

\paragraph{1181年（宋淳熙八年；公元换算据《宋史》本纪纪年）：南宋}
\noindent \eventanchor{SYD-1181-SS-207}{丙辰，以臨安疫，分命醫官診視軍民}\par

\paragraph{1182 年（共享年序桥接）：南宋、金与西夏（年序桥接）}
\noindent \eventanchor{SY-1182-CHRONOLOGY-BRIDGE}{【C级年序桥接】保留1182 年的多政权并行检索入口；本条不主张所列政权在当年发生直接互动，具体事件须另以单项材料入账。}\par

\paragraph{1182年（宋淳熙九年；公元换算据《宋史》本纪纪年）：南宋}
\noindent \eventanchor{SYD-1182-SS-056}{八月己亥朔，詔紹興民戶去歲已納夏稅應減者三十萬緡，理爲今年之數}\par

\paragraph{1182年（宋淳熙九年；公元换算据《宋史》本纪纪年）：南宋}
\noindent \eventanchor{SYD-1182-SS-208}{丙子，以子彤爲容州觀察使，封安定郡王}\par

\paragraph{1183 年（共享年序桥接）：南宋、金与西夏（年序桥接）}
\noindent \eventanchor{SY-1183-CHRONOLOGY-BRIDGE}{【C级年序桥接】保留1183 年的多政权并行检索入口；本条不主张所列政权在当年发生直接互动，具体事件须另以单项材料入账。}\par

\paragraph{1183年（宋淳熙十年；公元换算据《宋史》本纪纪年）：南宋}
从 《宋史》卷035〈淳熙十年〉 读 \eventanchor{SYD-1183-SS-057}{丁-\{丑\}-，詔除災傷州縣淳熙八年欠稅}，只能先确认 1183年（宋淳熙十年；公元换算据《宋史》本纪纪年） 的 南宋 与题名所示的具体事项。它并非抽象的年号标记：本纪年条记录政令或机构处置；实施背景须参照制度材料。；因此后面的选择受到 可在同年及后续政令中追踪；条文不单独证明地方执行。 的牵动。该判断不超出材料可承载范围，因为 《宋史》为元末官修，本纪保存宋廷纪年与处置；战果、数字、地方执行及对方动机须以编年、会要、对方史料或地方材料互校。

\paragraph{1183年（宋淳熙十年；公元换算据《宋史》本纪纪年）：南宋}
\noindent \eventanchor{SYD-1183-SS-209}{丁亥，金遣完顏婆盧火等來賀明年正旦}\par

\paragraph{1184 年（共享年序桥接）：南宋、金与西夏（年序桥接）}
\noindent \eventanchor{SY-1184-CHRONOLOGY-BRIDGE}{【C级年序桥接】保留1184 年的多政权并行检索入口；本条不主张所列政权在当年发生直接互动，具体事件须另以单项材料入账。}\par

\paragraph{1184年（宋淳熙十一年；公元换算据《宋史》本纪纪年）：南宋}
\noindent \eventanchor{SYD-1184-SS-058}{八月庚申，遣章森使金賀正旦}\par

\paragraph{1184年（宋淳熙十一年；公元换算据《宋史》本纪纪年）：南宋}
\noindent \eventanchor{SYD-1184-SS-210}{丙午，詔江東、西路諸監司，義役、差役從民便}\par

\paragraph{1185 年（共享年序桥接）：南宋、金与西夏（年序桥接）}
\noindent \eventanchor{SY-1185-CHRONOLOGY-BRIDGE}{【C级年序桥接】保留1185 年的多政权并行检索入口；本条不主张所列政权在当年发生直接互动，具体事件须另以单项材料入账。}\par

\paragraph{1185年（宋淳熙十二年；公元换算据《宋史》本纪纪年）：南宋}
\noindent \eventanchor{SYD-1185-SS-059}{八月癸亥，詔太上皇壽八十，令有司議慶壽禮}\par

\paragraph{1185年（宋淳熙十二年；公元换算据《宋史》本纪纪年）：南宋}
\noindent \eventanchor{SYD-1185-SS-211}{丙戌，詔恤潮州、台州被水之家}\par

\paragraph{1186 年（共享年序桥接）：南宋、金与西夏（年序桥接）}
\noindent \eventanchor{SY-1186-CHRONOLOGY-BRIDGE}{【C级年序桥接】保留1186 年的多政权并行检索入口；本条不主张所列政权在当年发生直接互动，具体事件须另以单项材料入账。}\par

\paragraph{1186年（宋淳熙十三年；公元换算据《宋史》本纪纪年）：南宋}
\noindent \eventanchor{SYD-1186-SS-060}{丙申，賜沖晦處士郭雍號曰頤正先生，仍遣官就問雍所欲言，備錄來上}\par

\paragraph{1186年（宋淳熙十三年；公元换算据《宋史》本纪纪年）：南宋}
\noindent \eventanchor{SYD-1186-SS-212}{丙寅，梁克家罷爲觀文殿大學士、醴泉觀使兼侍讀}\par

\paragraph{1187 年（共享年序桥接）：南宋、金与西夏（年序桥接）}
\noindent \eventanchor{SY-1187-CHRONOLOGY-BRIDGE}{【C级年序桥接】保留1187 年的多政权并行检索入口；本条不主张所列政权在当年发生直接互动，具体事件须另以单项材料入账。}\par

\paragraph{1187年（宋淳熙十四年；公元换算据《宋史》本纪纪年）：南宋}
\noindent \eventanchor{SYD-1187-SS-061}{八月辛未，賜度牒一百道、米四萬五千石，備振紹興府饑}\par

\paragraph{1187年（宋淳熙十四年；公元换算据《宋史》本纪纪年）：南宋}
\noindent \eventanchor{SYD-1187-SS-213}{丙辰，命臨安府捕蝗，募民輸米振濟}\par

\paragraph{1188 年（共享年序桥接）：南宋、金与西夏（年序桥接）}
\noindent \eventanchor{SY-1188-CHRONOLOGY-BRIDGE}{【C级年序桥接】保留1188 年的多政权并行检索入口；本条不主张所列政权在当年发生直接互动，具体事件须另以单项材料入账。}\par

\paragraph{1188年（宋淳熙十五年；公元换算据《宋史》本纪纪年）：南宋}
\noindent \eventanchor{SYD-1188-SS-062}{丁巳，詔修《高宗實錄》}\par

\paragraph{1188年（宋淳熙十五年；公元换算据《宋史》本纪纪年）：南宋}
从 《宋史》卷035〈淳熙十五年〉 读 \eventanchor{SYD-1188-SS-214}{己-\{丑\}-，詔減臨安、紹興府囚罪一等，釋杖以下，民緣欑宮役者蠲其賦}，只能先确认 1188年（宋淳熙十五年；公元换算据《宋史》本纪纪年） 的 南宋 与题名所示的具体事项。它并非抽象的年号标记：本纪年条记录政令或机构处置；实施背景须参照制度材料。；因此后面的选择受到 可在同年及后续政令中追踪；条文不单独证明地方执行。 的牵动。该判断不超出材料可承载范围，因为 《宋史》为元末官修，本纪保存宋廷纪年与处置；战果、数字、地方执行及对方动机须以编年、会要、对方史料或地方材料互校。

\paragraph{1189年：南宋}
\noindent \eventanchor{SS-1189-EVENT-128}{孝宗去世后光宗即位}\par

\paragraph{1189年（宋淳熙十六年；公元换算据《宋史》本纪纪年）：南宋}
\noindent \eventanchor{SYD-1189-SS-063}{甲午，封孫抦爲嘉國公}\par

\paragraph{1189年（宋淳熙十六年；公元换算据《宋史》本纪纪年）：南宋}
\noindent \eventanchor{SYD-1189-SS-215}{壬戌，下詔傳位皇太子}\par

\paragraph{1190 年（共享年序桥接）：南宋、金与西夏（年序桥接）}
\noindent \eventanchor{SY-1190-CHRONOLOGY-BRIDGE}{【C级年序桥接】保留1190 年的多政权并行检索入口；本条不主张所列政权在当年发生直接互动，具体事件须另以单项材料入账。}\par

\paragraph{1190年（宋紹熙元年；公元换算据《宋史》本纪纪年）：南宋}
\noindent \eventanchor{SYD-1190-SS-064}{八月辛卯，立任子中銓人吏部簾試法}\par

\paragraph{1190年（宋紹熙元年；公元换算据《宋史》本纪纪年）：南宋}
\noindent \eventanchor{SYD-1190-SS-216}{丙申，以上供等錢償廣州放免身丁錢數}\par

\paragraph{1191 年（共享年序桥接）：南宋、金与西夏（年序桥接）}
\noindent \eventanchor{SY-1191-CHRONOLOGY-BRIDGE}{【C级年序桥接】保留1191 年的多政权并行检索入口；本条不主张所列政权在当年发生直接互动，具体事件须另以单项材料入账。}\par

\paragraph{1191年（宋紹熙二年；公元换算据《宋史》本纪纪年）：南宋}
\noindent \eventanchor{SYD-1191-SS-065}{光宗紹熙二年，加諡受命中興全功至德聖神武文昭仁憲孝皇帝}\par

\paragraph{1191年（宋紹熙二年；公元换算据《宋史》本纪纪年）：南宋}
\noindent \eventanchor{SYD-1191-SS-217}{丙申，詔侍從、兩省、台諫及在外侍從之臣，各舉所知嘗任監司、郡守可充郎官、卿監及資歷未深可充諸職事官者各三人}\par

\paragraph{1192 年（共享年序桥接）：南宋、金与西夏（年序桥接）}
\noindent \eventanchor{SY-1192-CHRONOLOGY-BRIDGE}{【C级年序桥接】保留1192 年的多政权并行检索入口；本条不主张所列政权在当年发生直接互动，具体事件须另以单项材料入账。}\par

\paragraph{1192年（宋紹熙三年；公元换算据《宋史》本纪纪年）：南宋}
\noindent \eventanchor{SYD-1192-SS-066}{八月甲寅，詔兩淮行鐵錢交子}\par

\paragraph{1192年（宋紹熙三年；公元换算据《宋史》本纪纪年）：南宋}
\noindent \eventanchor{SYD-1192-SS-218}{兵部尚書羅點、給事中尤袤、中書舍人黃裳皆上疏請帝朝重華宮，吏部尚書趙汝愚亦因面對以請，帝開納}\par

\paragraph{1193年（宋紹熙四年；公元换算据《宋史》本纪纪年）：南宋}
\noindent \eventanchor{SYD-1193-SS-067}{明日會慶節，帝以疾不果朝，丞相葛邲率百官賀于重華宮}\par

\paragraph{1193年（宋紹熙四年；公元换算据《宋史》本纪纪年）：南宋}
\noindent \eventanchor{SYD-1193-SS-219}{丙寅，貸淮西民市牛錢}\par

\paragraph{1194年：南宋}
\noindent \eventanchor{SS-1194-EVENT-129}{宁宗即位，韩侂胄等在中枢取得影响}\par

\paragraph{1194年（宋紹熙五年；公元换算据《宋史》本纪纪年）：南宋}
\noindent \eventanchor{SYD-1194-SS-068}{丙辰，侍講黃裳、秘書少監孫逢吉等再上疏以請}\par

\paragraph{1194年（宋紹熙五年；公元换算据《宋史》本纪纪年）：南宋}
\noindent \eventanchor{SYD-1194-SS-220}{丙戌，權戶部侍郎袁說友入對，請朝重華宮}\par

\paragraph{1195 年（共享年序桥接）：南宋、金与西夏（年序桥接）}
\noindent \eventanchor{SY-1195-CHRONOLOGY-BRIDGE}{【C级年序桥接】保留1195 年的多政权并行检索入口；本条不主张所列政权在当年发生直接互动，具体事件须另以单项材料入账。}\par

\paragraph{1195年（宋慶元元年；公元换算据《宋史》本纪纪年）：南宋}
\noindent \eventanchor{SYD-1195-SS-069}{八月己巳，詔內外諸軍主帥條奏武備邊防之策以聞}\par

\paragraph{1195年（宋慶元元年；公元换算据《宋史》本纪纪年）：南宋}
核验 \eventanchor{SYD-1195-SS-221}{丙午，以監察-\{御\}-史胡紘言，責授趙汝愚寧遠軍節度副使、永州安置} 时，应回到 《宋史》卷037〈慶元元年〉 所保存的 1195年（宋慶元元年；公元换算据《宋史》本纪纪年） 记录，而不以事后概括代替它。南宋 围绕题名所示对象作出的处置有其条件：本纪年条记录政令或机构处置；实施背景须参照制度材料。；它把随后可见的选择推向 可在同年及后续政令中追踪；条文不单独证明地方执行。。《宋史》为元末官修，本纪保存宋廷纪年与处置；战果、数字、地方执行及对方动机须以编年、会要、对方史料或地方材料互校。

\paragraph{1196年：南宋}
\noindent \eventanchor{SS-1196-EVENT-130}{庆元党禁限制朱熹及相关士人的政治与学术活动}\par

\paragraph{1196年（宋慶元二年；公元换算据《宋史》本纪纪年）：南宋}
\noindent \eventanchor{SYD-1196-SS-070}{丙辰，乙太常少卿胡紘請，權住進擬偽學之黨}\par

\paragraph{1196年（宋慶元二年；公元换算据《宋史》本纪纪年）：南宋}
\noindent \eventanchor{SYD-1196-SS-222}{丙戌，減諸路死罪囚，釋流以下}\par

\paragraph{1197 年（共享年序桥接）：南宋、金与西夏（年序桥接）}
\noindent \eventanchor{SY-1197-CHRONOLOGY-BRIDGE}{【C级年序桥接】保留1197 年的多政权并行检索入口；本条不主张所列政权在当年发生直接互动，具体事件须另以单项材料入账。}\par

\paragraph{1197年（宋慶元三年；公元换算据《宋史》本纪纪年）：南宋}
\noindent \eventanchor{SYD-1197-SS-071}{八月戊子，復置嚴州神泉監}\par

\paragraph{1197年（宋慶元三年；公元换算据《宋史》本纪纪年）：南宋}
\noindent \eventanchor{SYD-1197-SS-223}{丙戌，金遣完顏愈來賀瑞慶節}\par

\paragraph{1198 年（共享年序桥接）：南宋、金与西夏（年序桥接）}
\noindent \eventanchor{SY-1198-CHRONOLOGY-BRIDGE}{【C级年序桥接】保留1198 年的多政权并行检索入口；本条不主张所列政权在当年发生直接互动，具体事件须另以单项材料入账。}\par

\paragraph{1198年（宋慶元四年；公元换算据《宋史》本纪纪年）：南宋}
\noindent \eventanchor{SYD-1198-SS-072}{八月丁卯朔，以久雨，決系囚}\par

\paragraph{1198年（宋慶元四年；公元换算据《宋史》本纪纪年）：南宋}
\noindent \eventanchor{SYD-1198-SS-224}{丙戌，詔以太上皇聖躬清復，率群臣上壽}\par

\paragraph{1199 年（共享年序桥接）：南宋、金与西夏（年序桥接）}
\noindent \eventanchor{SY-1199-CHRONOLOGY-BRIDGE}{【C级年序桥接】保留1199 年的多政权并行检索入口；本条不主张所列政权在当年发生直接互动，具体事件须另以单项材料入账。}\par

\paragraph{1199年（宋慶元五年；公元换算据《宋史》本纪纪年）：南宋}
\noindent \eventanchor{SYD-1199-SS-073}{丙辰，遣朱致知使金賀正旦}\par

\paragraph{1199年（宋慶元五年；公元换算据《宋史》本纪纪年）：南宋}
\noindent \eventanchor{SYD-1199-SS-225}{丙戌，詔減諸路流囚，釋杖以下，推恩如慶壽故事}\par

\paragraph{1200 年（共享年序桥接）：南宋、金与西夏（年序桥接）}
\noindent \eventanchor{SY-1200-CHRONOLOGY-BRIDGE}{【C级年序桥接】保留1200 年的多政权并行检索入口；本条不主张所列政权在当年发生直接互动，具体事件须另以单项材料入账。}\par

\paragraph{1200年（宋慶元六年；公元换算据《宋史》本纪纪年）：南宋}
\noindent \eventanchor{SYD-1200-SS-074}{丙寅，詔大理、三衙、臨安府及諸路闕雨州縣釋杖以下囚}\par

\paragraph{1200年（宋慶元六年；公元换算据《宋史》本纪纪年）：南宋}
1200年（宋慶元六年；公元换算据《宋史》本纪纪年），南宋 的 \eventanchor{SYD-1200-SS-226}{丁-\{丑\}-，詔三省、樞密院擇臣僚封事可行者以聞} 见于 《宋史》卷037〈慶元六年〉。题名留下的行动、文书、城路或人名构成此处唯一不应被抹平的材料细节；本纪年条记录政令或机构处置；实施背景须参照制度材料。。其影响只能谨慎写到 可在同年及后续政令中追踪；条文不单独证明地方执行。。《宋史》为元末官修，本纪保存宋廷纪年与处置；战果、数字、地方执行及对方动机须以编年、会要、对方史料或地方材料互校。

\paragraph{1201 年（共享年序桥接）：南宋、金与西夏（年序桥接）}
\noindent \eventanchor{SY-1201-CHRONOLOGY-BRIDGE}{【C级年序桥接】保留1201 年的多政权并行检索入口；本条不主张所列政权在当年发生直接互动，具体事件须另以单项材料入账。}\par

\paragraph{1201年（宋嘉泰元年；公元换算据《宋史》本纪纪年）：南宋}
\noindent \eventanchor{SYD-1201-SS-075}{頒《慶元寬恤詔令》、《役法撮要》}\par

\paragraph{1201年（宋嘉泰元年；公元换算据《宋史》本纪纪年）：南宋}
\noindent \eventanchor{SYD-1201-SS-227}{丙午，詔文武臣無寓居州任厘務官，著爲令}\par

\paragraph{1202 年（共享年序桥接）：南宋、金与西夏（年序桥接）}
\noindent \eventanchor{SY-1202-CHRONOLOGY-BRIDGE}{【C级年序桥接】保留1202 年的多政权并行检索入口；本条不主张所列政权在当年发生直接互动，具体事件须另以单项材料入账。}\par

\paragraph{1202年（宋嘉泰二年；公元换算据《宋史》本纪纪年）：南宋}
\noindent \eventanchor{SYD-1202-SS-076}{丙寅，嗣秀王伯圭薨，追封崇王，諡曰憲靖}\par

\paragraph{1202年（宋嘉泰二年；公元换算据《宋史》本纪纪年）：南宋}
\noindent \eventanchor{SYD-1202-SS-228}{丁亥，修《高宗正史》、《寶訓》}\par

\paragraph{1203 年（共享年序桥接）：南宋、金与西夏（年序桥接）}
\noindent \eventanchor{SY-1203-CHRONOLOGY-BRIDGE}{【C级年序桥接】保留1203 年的多政权并行检索入口；本条不主张所列政权在当年发生直接互动，具体事件须另以单项材料入账。}\par

\paragraph{1203年（宋嘉泰三年；公元换算据《宋史》本纪纪年）：南宋}
\noindent \eventanchor{SYD-1203-SS-077}{八月壬寅，增置襄陽騎軍}\par

\paragraph{1203年（宋嘉泰三年；公元换算据《宋史》本纪纪年）：南宋}
\noindent \eventanchor{SYD-1203-SS-229}{丙午，出封樁庫兩淮交子一百萬，命轉運司收民間鐵錢}\par

\paragraph{1204年：南宋与金}
\noindent \eventanchor{SS-1204-EVENT-131}{韩侂胄推动北伐准备}\par

\paragraph{1204年（宋嘉泰四年；公元换算据《宋史》本纪纪年）：南宋}
\noindent \eventanchor{SYD-1204-SS-078}{丙辰，除靜江府、昭州折布錢}\par

\paragraph{1204年（宋嘉泰四年；公元换算据《宋史》本纪纪年）：南宋}
\noindent \eventanchor{SYD-1204-SS-230}{丙申，置諸軍帳前雄校，以軍官子孫補之}\par

\paragraph{1205 年（共享年序桥接）：南宋、金与西夏（年序桥接）}
\noindent \eventanchor{SY-1205-CHRONOLOGY-BRIDGE}{【C级年序桥接】保留1205 年的多政权并行检索入口；本条不主张所列政权在当年发生直接互动，具体事件须另以单项材料入账。}\par

\paragraph{1205年（宋開禧元年；公元换算据《宋史》本纪纪年）：南宋}
\noindent \eventanchor{SYD-1205-SS-079}{八月丙戌朔，蠲兩浙闕雨州縣贓賞錢}\par

\paragraph{1205年（宋開禧元年；公元换算据《宋史》本纪纪年）：南宋}
\noindent \eventanchor{SYD-1205-SS-231}{丙寅，升嘉定府爲嘉慶軍}\par

\paragraph{1206年（宋开禧二年、金泰和六年；干支、月日待核；公元换算据双方本纪年号对照）：南宋与金}
\noindent \eventanchor{SS-1206-EVENT-132}{南宋发动开禧北伐，战事不利}\par

\paragraph{1206年（宋開禧二年；公元换算据《宋史》本纪纪年）：南宋}
\noindent \eventanchor{SYD-1206-SS-080}{八月丙寅，有司上《開禧刑名斷例》}\par

\paragraph{1206年（宋開禧二年；公元换算据《宋史》本纪纪年）：南宋}
\noindent \eventanchor{SYD-1206-SS-232}{罷旱傷州軍比較租賦一年}\par

\paragraph{1207年：南宋}
\noindent \eventanchor{SS-1207-EVENT-133}{韩侂胄被杀，朝廷转向结束北伐}\par

\paragraph{1207年（宋開禧三年；公元换算据《宋史》本纪纪年）：南宋}
\noindent \eventanchor{SYD-1207-SS-081}{辛巳，再奪鄧友龍五官、南雄州安置，尋除名，徙循州}\par

\paragraph{1207年（宋開禧三年；公元换算据《宋史》本纪纪年）：南宋}
\noindent \eventanchor{SYD-1207-SS-233}{丙戌，命淮西轉運司措置雄淮軍}\par

\paragraph{1208年（宋嘉定元年、金泰和八年；干支、月日待核；公元换算据双方本纪年号对照）：南宋与金}
\noindent \eventanchor{SS-1208-EVENT-134}{南宋与金订立嘉定和议}\par

\paragraph{1208年（宋嘉定元年；公元换算据《宋史》本纪纪年）：南宋}
\noindent \eventanchor{SYD-1208-SS-082}{八月戊辰朔，發米振貧民}\par

\paragraph{1208年（宋嘉定元年；公元换算据《宋史》本纪纪年）：南宋}
\noindent \eventanchor{SYD-1208-SS-234}{丙申，幸太乙宮、明慶寺禱雨}\par

\paragraph{1209年（宋嘉定二年；公元换算据《宋史》本纪纪年）：南宋}
\noindent \eventanchor{SYD-1209-SS-083}{丙戌，發米十萬石振兩淮饑民}\par

\paragraph{1209年（宋嘉定二年；公元换算据《宋史》本纪纪年）：南宋}
\noindent \eventanchor{SYD-1209-SS-235}{丙戌，金遣使來賀明年正旦}\par

\paragraph{1210年（宋嘉定三年；公元换算据《宋史》本纪纪年）：南宋}
\noindent \eventanchor{SYD-1210-SS-084}{丙辰，以久雨，釋兩浙州縣繫囚}\par

\paragraph{1210年（宋嘉定三年；公元换算据《宋史》本纪纪年）：南宋}
\noindent \eventanchor{SYD-1210-SS-236}{丙寅，湖南賊羅世傳縛李元礪以降，峒寇悉平}\par

\paragraph{1211年（宋嘉定四年；公元换算据《宋史》本纪纪年）：南宋}
\noindent \eventanchor{SYD-1211-SS-085}{丙午，賜黑風峒名曰效忠}\par

\paragraph{1211年（宋嘉定四年；公元换算据《宋史》本纪纪年）：南宋}
\noindent \eventanchor{SYD-1211-SS-237}{丙午，詔湖南、江西諸州經賊蹂踐者，監司、守臣考縣令安集之實，第其能否以聞}\par

\paragraph{1212年（宋嘉定五年；公元换算据《宋史》本纪纪年）：南宋}
《宋史》卷039〈嘉定五年〉 为 \eventanchor{SYD-1212-SS-086}{八月甲戌朔，-\{御\}-後殿，復膳} 提供年序凭据。南宋 在 1212年（宋嘉定五年；公元换算据《宋史》本纪纪年） 面对的不是没有空间的政权名，而是题名所列的事务及其可调度的对象；本纪年条记录政令或机构处置；实施背景须参照制度材料。。它使后续叙事必须继续追问 可在同年及后续政令中追踪；条文不单独证明地方执行。。与此同时，《宋史》为元末官修，本纪保存宋廷纪年与处置；战果、数字、地方执行及对方动机须以编年、会要、对方史料或地方材料互校。

\paragraph{1212年（宋嘉定五年；公元换算据《宋史》本纪纪年）：南宋}
\noindent \eventanchor{SYD-1212-SS-238}{冬十月辛巳，詔諸路總領官歲舉堪將帥者二三人，安撫、提刑舉可備將材者各二人}\par

\paragraph{1213年（宋嘉定六年；公元换算据《宋史》本纪纪年）：南宋}
\noindent \eventanchor{SYD-1213-SS-087}{八月己巳朔，詔諸路監司、帥臣舉所部官吏之才行卓絕、績用章著者}\par

\paragraph{1213年（宋嘉定六年；公元换算据《宋史》本纪纪年）：南宋}
\noindent \eventanchor{SYD-1213-SS-239}{丙申，以雷發非時，下罪己詔}\par

\paragraph{1214年（宋嘉定七年；公元换算据《宋史》本纪纪年）：南宋}
\noindent \eventanchor{SYD-1214-SS-088}{八月癸巳朔，罷關外四州所增方田稅}\par

\paragraph{1214年（宋嘉定七年；公元换算据《宋史》本纪纪年）：南宋}
\noindent \eventanchor{SYD-1214-SS-240}{罷四川制置大使司所開鹽井}\par

\paragraph{1215年（宋嘉定八年；公元换算据《宋史》本纪纪年）：南宋}
\noindent \eventanchor{SYD-1215-SS-089}{八年春正月辛未，命師禹嗣秀王}\par

\paragraph{1215年（宋嘉定八年；公元换算据《宋史》本纪纪年）：南宋}
\eventanchor{SYD-1215-SS-241}{八月己-\{丑\}-，賜張栻諡曰宣} 所指的 1215年（宋嘉定八年；公元换算据《宋史》本纪纪年） 节点，保留了 南宋 处理题名所涉人事、资源、地点或制度的痕迹；可复核的出处为 《宋史》卷039〈嘉定八年〉。前一层条件不是宿命，而是 本纪年条提供日期和中枢可见行动；前因不据单条补定。；随之可追踪的结果为 可回查同年及后续条目，地方影响仍须另证。。《宋史》为元末官修，本纪保存宋廷纪年与处置；战果、数字、地方执行及对方动机须以编年、会要、对方史料或地方材料互校。

\paragraph{1216年（宋嘉定九年；公元换算据《宋史》本纪纪年）：南宋}
\noindent \eventanchor{SYD-1216-SS-090}{丙寅，金遣使來賀瑞慶節}\par

\paragraph{1216年（宋嘉定九年；公元换算据《宋史》本纪纪年）：南宋}
\noindent \eventanchor{SYD-1216-SS-242}{丙子，命諸州招填軍籍}\par

\paragraph{1217年：南宋与金}
\noindent \eventanchor{SS-1217-EVENT-135}{金军南攻四川、两淮等方向}\par

\paragraph{1217年（宋嘉定十年；公元换算据《宋史》本纪纪年）：南宋}
\noindent \eventanchor{SYD-1217-SS-091}{八月乙丑，詔監司、郡守各舉威勇才略可將帥者二人}\par

\paragraph{1217年（宋嘉定十年；公元换算据《宋史》本纪纪年）：南宋}
\noindent \eventanchor{SYD-1217-SS-243}{丙辰，詔江淮制置使李玨、京湖制置使趙方措置調遣，仍聽便宜行事}\par

\paragraph{1218年（宋嘉定十一年；公元换算据《宋史》本纪纪年）：南宋}
\noindent \eventanchor{SYD-1218-SS-092}{丁亥，詔侍從、臺諫、兩省官集議平戎、禦戎、和戎三策}\par

\paragraph{1218年（宋嘉定十一年；公元换算据《宋史》本纪纪年）：南宋}
\noindent \eventanchor{SYD-1218-SS-244}{丁酉，詔四川忠義人立功，賞視官軍}\par

\paragraph{1219年（宋嘉定十二年；公元换算据《宋史》本纪纪年）：南宋}
\noindent \eventanchor{SYD-1219-SS-093}{八月戊辰，復合利州東、西路爲一}\par

\paragraph{1219年（宋嘉定十二年；公元换算据《宋史》本纪纪年）：南宋}
《宋史》卷040〈嘉定十二年〉 为 \eventanchor{SYD-1219-SS-245}{丁-\{丑\}-，雅州蠻入盧山縣} 提供年序凭据。南宋 在 1219年（宋嘉定十二年；公元换算据《宋史》本纪纪年） 面对的不是没有空间的政权名，而是题名所列的事务及其可调度的对象；本纪年条提供日期和中枢可见行动；前因不据单条补定。。它使后续叙事必须继续追问 可回查同年及后续条目，地方影响仍须另证。。与此同时，《宋史》为元末官修，本纪保存宋廷纪年与处置；战果、数字、地方执行及对方动机须以编年、会要、对方史料或地方材料互校。

\paragraph{1220年（宋嘉定十三年；公元换算据《宋史》本纪纪年）：南宋}
\noindent \eventanchor{SYD-1220-SS-094}{丙辰，四川宣撫司招黎人土丁，降之}\par

\paragraph{1220年（宋嘉定十三年；公元换算据《宋史》本纪纪年）：南宋}
\noindent \eventanchor{SYD-1220-SS-246}{丁巳，黎州土丁叛，遣兵討之}\par

\paragraph{1221年（宋嘉定十四年；公元换算据《宋史》本纪纪年）：南宋}
\noindent \eventanchor{SYD-1221-SS-095}{八月乙卯，賜史彌遠家廟}\par

\paragraph{1221年（宋嘉定十四年；公元换算据《宋史》本纪纪年）：南宋}
\noindent \eventanchor{SYD-1221-SS-247}{丙寅，夏人復以書來四川趣會兵}\par

\paragraph{1222年（宋嘉定十五年；公元换算据《宋史》本纪纪年）：南宋}
\noindent \eventanchor{SYD-1222-SS-096}{八月己卯，命戶部詳議義役}\par

\paragraph{1222年（宋嘉定十五年；公元换算据《宋史》本纪纪年）：南宋}
\noindent \eventanchor{SYD-1222-SS-248}{丁巳，復以隨州三關隸德安府，置關使}\par

\paragraph{1223年（宋嘉定十六年；公元换算据《宋史》本纪纪年）：南宋}
\noindent \eventanchor{SYD-1223-SS-097}{丁卯，以道州民饑，詔發米振之}\par

\paragraph{1223年（宋嘉定十六年；公元换算据《宋史》本纪纪年）：南宋}
\noindent \eventanchor{SYD-1223-SS-249}{冬十一月辛亥，以太平州大水，詔振恤之}\par

\paragraph{1224年：南宋}
\noindent \eventanchor{SS-1224-EVENT-136}{宁宗去世，理宗即位，史弥远继续主导中枢}\par

\paragraph{1224年（宋嘉定十七年；公元换算据《宋史》本纪纪年）：南宋}
\noindent \eventanchor{SYD-1224-SS-098}{八月乙亥，罷通州天賜鹽場}\par

\paragraph{1224年（宋嘉定十七年；公元换算据《宋史》本纪纪年）：南宋}
\noindent \eventanchor{SYD-1224-SS-250}{癸亥，命淮東西、湖北路轉運司提督營屯田}\par

\paragraph{1225年（宋寶慶元年；公元换算据《宋史》本纪纪年）：南宋}
\noindent \eventanchor{SYD-1225-SS-099}{丙寅，詔不熄爲保康軍承宣使、嗣濮王}\par

\paragraph{1225年（宋寶慶元年；公元换算据《宋史》本纪纪年）：南宋}
\noindent \eventanchor{SYD-1225-SS-251}{壬寅，帝兩請皇太后垂簾，不允}\par

\paragraph{1226年（宋寶慶二年；公元换算据《宋史》本纪纪年）：南宋}
\noindent \eventanchor{SYD-1226-SS-100}{八月乙巳，濟王竑追降巴陵縣公}\par

\paragraph{1226年（宋寶慶二年；公元换算据《宋史》本纪纪年）：南宋}
\eventanchor{SYD-1226-SS-252}{冬十月甲申，詔《寧宗-\{御\}-集》閣以「寶章」爲名，仍置學士、待制員} 所指的 1226年（宋寶慶二年；公元换算据《宋史》本纪纪年） 节点，保留了 南宋 处理题名所涉人事、资源、地点或制度的痕迹；可复核的出处为 《宋史》卷041〈寶慶二年〉。前一层条件不是宿命，而是 本纪年条记录政令或机构处置；实施背景须参照制度材料。；随之可追踪的结果为 可在同年及后续政令中追踪；条文不单独证明地方执行。。《宋史》为元末官修，本纪保存宋廷纪年与处置；战果、数字、地方执行及对方动机须以编年、会要、对方史料或地方材料互校。

\paragraph{1227年（宋寶慶三年；公元换算据《宋史》本纪纪年）：南宋}
\noindent \eventanchor{SYD-1227-SS-101}{己巳，詔：「朕觀朱熹集注《大學》、《論語》、《孟子》、《中庸》，發揮聖賢蘊奧，有補治道，朕勵志講學，緬懷典刑，可特贈熹太師，追封信國公}\par

\paragraph{1227年（宋寶慶三年；公元换算据《宋史》本纪纪年）：南宋}
\noindent \eventanchor{SYD-1227-SS-253}{八月庚戌，詔謝氏特封通義郡夫人}\par

\paragraph{1228年（宋紹定元年；公元换算据《宋史》本纪纪年）：南宋}
\noindent \eventanchor{SYD-1228-SS-102}{丁酉，詔申嚴皇城司給符之制，照闌入法}\par

\paragraph{1229年（宋紹定二年；公元换算据《宋史》本纪纪年）：南宋}
\noindent \eventanchor{SYD-1229-SS-103}{冬十月壬戌，詔台州水災，除民田租及茶、鹽、酒酤諸雜稅，郡縣抑納者監司察之}\par

\paragraph{1229年（宋紹定二年；公元换算据《宋史》本纪纪年）：南宋}
\noindent \eventanchor{SYD-1229-SS-254}{二月庚戌，詔歲舉廉吏或犯奸贓，保任同坐，監司、守臣其申嚴覺察}\par

\paragraph{1230年（宋紹定三年；公元换算据《宋史》本纪纪年）：南宋}
\noindent \eventanchor{SYD-1230-SS-104}{丁卯，冊命貴妃謝氏爲皇后}\par

\paragraph{1230年（宋紹定三年；公元换算据《宋史》本纪纪年）：南宋}
\noindent \eventanchor{SYD-1230-SS-255}{乙丑，詔免明年元會禮}\par

\paragraph{1231年（宋紹定四年；公元换算据《宋史》本纪纪年）：南宋}
\noindent \eventanchor{SYD-1231-SS-105}{八月己未，大元兵破武休，入興元，攻仙人關}\par

\paragraph{1231年（宋紹定四年；公元换算据《宋史》本纪纪年）：南宋}
\noindent \eventanchor{SYD-1231-SS-256}{丙子，詔起復孟珙從義郎、京西路分，棗陽軍駐紮}\par

\paragraph{1232年（宋紹定五年；公元换算据《宋史》本纪纪年）：南宋}
\noindent \eventanchor{SYD-1232-SS-106}{壬辰，史嵩之進大理卿、權刑部侍郎、京湖安撫制置使、知襄陽府}\par

\paragraph{1232年（宋紹定五年；公元换算据《宋史》本纪纪年）：南宋}
\noindent \eventanchor{SYD-1232-SS-257}{比聞蘄州進士馮傑，本儒家，都大坑冶司抑爲爐戶，誅求日增，傑妻以憂死，其女繼之，弟大聲因赴訴，死于道路，傑知不免，毒其二子一妾，舉火自經而死}\par

\paragraph{1233年：南宋与蒙古}
\noindent \eventanchor{SS-1233-EVENT-137}{南宋与蒙古在攻金问题上形成短暂军事协作}\par

\paragraph{1233年（宋紹定六年；公元换算据《宋史》本纪纪年）：南宋}
1233年（宋紹定六年；公元换算据《宋史》本纪纪年） 的 《宋史》卷041〈紹定六年〉 记下 \eventanchor{SYD-1233-SS-107}{己-\{丑\}-，詔崔與之、李、鄭性之赴闕}。对 南宋 而言，题名所点出的对象、机构或行动是这条记录能够落地的线索；本纪年条记录政令或机构处置；实施背景须参照制度材料。。它在后续年序中所打开或加重的是：可在同年及后续政令中追踪；条文不单独证明地方执行。。《宋史》为元末官修，本纪保存宋廷纪年与处置；战果、数字、地方执行及对方动机须以编年、会要、对方史料或地方材料互校。

\paragraph{1233年（宋紹定六年；公元换算据《宋史》本纪纪年）：南宋}
\noindent \eventanchor{SYD-1233-SS-258}{詔袁韶奪職、罷祠祿}\par

\paragraph{1234年（宋端平元年、金天兴三年、蒙古窝阔台六年；干支、月日待核；公元换算据各方本纪年号对照）：南宋金蒙古}
\noindent \eventanchor{SS-1234-EVENT-138}{蔡州陷落后，南宋试图进入河南故地}\par

\paragraph{1234年（宋端平元年；公元换算据《宋史》本纪纪年）：南宋}
\noindent \eventanchor{SYD-1234-SS-108}{甲戌，朱揚祖、林拓朝謁八陵回，以圖進，上問諸陵相去幾何及陵前澗水新復，揚祖悉以對，上忍涕太息}\par

\paragraph{1234年（宋端平元年；公元换算据《宋史》本纪纪年）：南宋}
\noindent \eventanchor{SYD-1234-SS-259}{建陽縣盜發，衆數千人，焚劫邵武、麻沙、長平}\par

\paragraph{1235年：南宋与蒙古}
\noindent \eventanchor{SS-1235-EVENT-139}{蒙古军开始多方向进攻南宋边境}\par

\paragraph{1235年（宋端平二年；公元换算据《宋史》本纪纪年）：南宋}
\noindent \eventanchor{SYD-1235-SS-109}{丙辰，詔主管侍衛馬軍孟珙黃州駐紮，措置邊防}\par

\paragraph{1235年（宋端平二年；公元换算据《宋史》本纪纪年）：南宋}
\noindent \eventanchor{SYD-1235-SS-260}{賜進士吳叔告以下四百五十四人及第、出身有差}\par

\paragraph{1236年（宋端平三年；公元换算据《宋史》本纪纪年）：南宋}
\noindent \eventanchor{SYD-1236-SS-110}{甲午，詔：「沿江制置使陳韡應援淮東，授淮西制置使兼沿江制置副使史嵩之應援江陵、峽州江面上流}\par

\paragraph{1237年（宋嘉熙元年；公元换算据《宋史》本纪纪年）：南宋}
\noindent \eventanchor{SYD-1237-SS-111}{詔曾伯等十一人各官一轉}\par

\paragraph{1238年（宋嘉熙二年；公元换算据《宋史》本纪纪年）：南宋}
\noindent \eventanchor{SYD-1238-SS-112}{戊戌，詔：「近覽李奏，知蜀漸次收復，然創殘之餘，綏撫爲急，宜施蕩宥之澤}\par

\paragraph{1239年（宋嘉熙三年；公元换算据《宋史》本纪纪年）：南宋}
\noindent \eventanchor{SYD-1239-SS-113}{秉義郎李良守鄂州長壽縣，沒于戰陣，詔贈官三轉}\par

\paragraph{1240年（宋嘉熙四年；公元换算据《宋史》本纪纪年）：南宋}
\noindent \eventanchor{SYD-1240-SS-114}{又言：「張祥有保全趙彥呐、楊恢兩制置之功，敵人憚其果毅，宜見錄用}\par

\paragraph{1241年（宋淳祐元年；公元换算据《宋史》本纪纪年）：南宋}
\noindent \eventanchor{SYD-1241-SS-115}{尋以王安石謂「天命不足畏，祖宗不足法，人言不足恤」，爲萬世罪人，豈宜從祀孔子廟庭，黜之}\par

\paragraph{1242年（宋淳祐二年；公元换算据《宋史》本纪纪年）：南宋}
\noindent \eventanchor{SYD-1242-SS-116}{八月丁卯，詔淮東先鋒馬軍鄧淳、李海等揚州撻扒店之戰，宣勞居多，各官兩轉，餘推恩有差}\par

\paragraph{1243年（宋淳祐三年；公元换算据《宋史》本纪纪年）：南宋}
\noindent \eventanchor{SYD-1243-SS-117}{丙辰，安豐軍統領陳友直以王家堈戰功，與官兩轉，壬申，布衣王與之進所著《周禮訂議》，補下州文學}\par

\paragraph{1244年（宋淳祐四年；公元换算据《宋史》本纪纪年）：南宋}
\noindent \eventanchor{SYD-1244-SS-118}{詔<u>玠</u>趣上立功將士姓名等第，即與推恩}\par

\paragraph{1245年（宋淳祐五年；公元换算据《宋史》本纪纪年）：南宋}
\noindent \eventanchor{SYD-1245-SS-119}{詔王雲贈三秩，仍官其二子爲承信郎}\par

\paragraph{1246年（宋淳祐六年；公元换算据《宋史》本纪纪年）：南宋}
\noindent \eventanchor{SYD-1246-SS-120}{詔倪政贈官三轉，官一子承信郎}\par

\paragraph{1247年（宋淳祐七年；公元换算据《宋史》本纪纪年）：南宋}
\noindent \eventanchor{SYD-1247-SS-121}{丙申，以旱，避殿減膳}\par

\paragraph{1248年（宋淳祐八年；公元换算据《宋史》本纪纪年）：南宋}
\noindent \eventanchor{SYD-1248-SS-122}{詔命詞、給告身付之}\par

\paragraph{1249年（宋淳祐九年；公元换算据《宋史》本纪纪年）：南宋}
\noindent \eventanchor{SYD-1249-SS-123}{癸亥，詔給官田五百畝，命臨安府創慈幼局，收養道路遺棄初生嬰兒，仍置藥局療貧民疾病}\par

\paragraph{1250年（宋淳祐十年；公元换算据《宋史》本纪纪年）：南宋}
\noindent \eventanchor{SYD-1250-SS-124}{參知政事謝方叔、吳潛、簽書樞密院事徐清叟並乞解機政，詔不允}\par

\paragraph{1251年（宋淳祐十一年；公元换算据《宋史》本纪纪年）：南宋}
\noindent \eventanchor{SYD-1251-SS-125}{丁未，命呂文福廬州駐紮御前諸軍都統制}\par

\paragraph{1252年（宋淳祐十二年；公元换算据《宋史》本纪纪年）：南宋}
\noindent \eventanchor{SYD-1252-SS-126}{八月己未，詔來年省試仍舊用二月一日，殿試用四月十五日以前，庶免滯留遠方士子}\par

\paragraph{1253年：南宋与蒙古}
\noindent \eventanchor{SS-1253-EVENT-140}{蒙古征服大理，南宋西南战略环境改变}\par

\paragraph{1253年（宋寶祐元年；公元换算据《宋史》本纪纪年）：南宋}
\noindent \eventanchor{SYD-1253-SS-127}{寶祐元年春正月庚寅朔，詔以藝祖嫡系十一世孫嗣榮王與芮之子建安郡王孜爲皇子，改賜名禥，授崇慶軍節度使，進封永嘉郡王}\par

\paragraph{1254年（宋寶祐二年；公元换算据《宋史》本纪纪年）：南宋}
\noindent \eventanchor{SYD-1254-SS-128}{詔蒲擇之暫權四川制置司事}\par

\paragraph{1255年（宋寶祐三年；公元换算据《宋史》本纪纪年）：南宋}
《宋史》卷044〈寶祐三年〉 为 \eventanchor{SYD-1255-SS-129}{丙辰，謝方叔、徐清叟以-\{御\}-史朱應元言罷} 提供年序凭据。南宋 在 1255年（宋寶祐三年；公元换算据《宋史》本纪纪年） 面对的不是没有空间的政权名，而是题名所列的事务及其可调度的对象；本纪年条记录政令或机构处置；实施背景须参照制度材料。。它使后续叙事必须继续追问 可在同年及后续政令中追踪；条文不单独证明地方执行。。与此同时，《宋史》为元末官修，本纪保存宋廷纪年与处置；战果、数字、地方执行及对方动机须以编年、会要、对方史料或地方材料互校。

\paragraph{1256年（宋寶祐四年；公元换算据《宋史》本纪纪年）：南宋}
1256年（宋寶祐四年；公元换算据《宋史》本纪纪年），南宋 的 \eventanchor{SYD-1256-SS-130}{乙巳，以監察-\{御\}-史吳衍、翁應弼劾太學武學生劉黻等八人不率，詔拘管江西、湖南州軍，宗學生與亻匈等七人並削籍，拘管外宗正司} 见于 《宋史》卷044〈寶祐四年〉。题名留下的行动、文书、城路或人名构成此处唯一不应被抹平的材料细节；本纪年条记录政令或机构处置；实施背景须参照制度材料。。其影响只能谨慎写到 可在同年及后续政令中追踪；条文不单独证明地方执行。。《宋史》为元末官修，本纪保存宋廷纪年与处置；战果、数字、地方执行及对方动机须以编年、会要、对方史料或地方材料互校。

\paragraph{1257年（宋寶祐五年；公元换算据《宋史》本纪纪年）：南宋}
\noindent \eventanchor{SYD-1257-SS-131}{丙午，禁奸民作白衣會，監司、郡縣官等失覺察者坐罪}\par

\paragraph{1258年：南宋与蒙古}
\noindent \eventanchor{SS-1258-EVENT-141}{蒙哥发动大规模攻宋计划}\par

\paragraph{1258年（宋寶祐六年；公元换算据《宋史》本纪纪年）：南宋}
\noindent \eventanchor{SYD-1258-SS-132}{丙辰，給事中張鎮言：徐敏子曩帥廣右，嗜殺黷貨，流毒桂府}\par

\paragraph{1259年：南宋与蒙古}
\noindent \eventanchor{SS-1259-EVENT-142}{钓鱼城战事与蒙哥去世改变战争进程}\par

\paragraph{1259年（宋開慶元年；公元换算据《宋史》本纪纪年）：南宋}
\noindent \eventanchor{SYD-1259-SS-133}{庚辰，以趙與爲觀文殿學士、兩淮安撫制置使兼知揚州}\par

\paragraph{1260年（宋景定元年；公元换算据《宋史》本纪纪年）：南宋}
\noindent \eventanchor{SYD-1260-SS-134}{賈似道奏：「和出彼謀，豈容一切輕徇？倘以交鄰國之道來，當令入見}\par

\paragraph{1261年（宋景定二年；公元换算据《宋史》本纪纪年）：南宋}
\noindent \eventanchor{SYD-1261-SS-135}{’請宣付史館，永爲世程法}\par

\paragraph{1263年（宋景定四年；公元换算据《宋史》本纪纪年）：南宋}
\noindent \eventanchor{SYD-1263-SS-136}{八月甲寅，董宋臣以病乞收回恩命，請祠，詔賜告五月}\par

\paragraph{1264年：南宋}
\noindent \eventanchor{SS-1264-EVENT-143}{理宗去世，度宗即位，贾似道权势上升}\par

\paragraph{1264年（宋景定五年；公元换算据《宋史》本纪纪年）：南宋}
\eventanchor{SYD-1264-SS-137}{丁-\{丑\}-，詔避殿減膳，應中外臣僚許直言朝政闕失} 的可见载体是 《宋史》卷045〈景定五年〉；它把 南宋 放在 1264年（宋景定五年；公元换算据《宋史》本纪纪年） 的这一处置里。本纪年条记录政令或机构处置；实施背景须参照制度材料。 所以这不是可由相邻年份自动推出的结果；其后留下的条件是：可在同年及后续政令中追踪；条文不单独证明地方执行。。材料边界也须保留：《宋史》为元末官修，本纪保存宋廷纪年与处置；战果、数字、地方执行及对方动机须以编年、会要、对方史料或地方材料互校。

\paragraph{1265年（宋咸淳元年；公元换算据《宋史》本纪纪年）：南宋}
\noindent \eventanchor{SYD-1265-SS-138}{八月庚辰，命陳奕沿江按閱軍防，賜錢二十萬給用}\par

\paragraph{1266年（宋咸淳二年；公元换算据《宋史》本纪纪年）：南宋}
\noindent \eventanchor{SYD-1266-SS-139}{詔自咸淳三年爲始罷之}\par

\paragraph{1267年（宋咸淳三年；公元换算据《宋史》本纪纪年）：南宋}
\noindent \eventanchor{SYD-1267-SS-140}{八月辛酉，遣步帥陳奕率馬軍舟師巡邏江防}\par

\paragraph{1268年：南宋与蒙古}
\noindent \eventanchor{SS-1268-EVENT-144}{蒙古军开始围攻襄阳、樊城}\par

\paragraph{1268年（宋咸淳四年；公元换算据《宋史》本纪纪年）：南宋}
\noindent \eventanchor{SYD-1268-SS-141}{於是盧鉞等相繼論列方叔昨蜀、廣敗事，誤國殄民，今又違制擅進，削一秩罰輕}\par

\paragraph{1269年（宋咸淳五年；公元换算据《宋史》本纪纪年）：南宋}
\noindent \eventanchor{SYD-1269-SS-142}{八月戊寅，詔郡縣收民田租，毋巧計取贏，毋厚直折納，轉運司申嚴按劾}\par

\paragraph{1270年（宋咸淳六年；公元换算据《宋史》本纪纪年）：南宋}
\noindent \eventanchor{SYD-1270-SS-143}{癸亥，詔：「贛、吉、南安境數被寇，雖有砦卒，寇出沒無時，莫能相救}\par

\paragraph{1271年（宋咸淳七年；公元换算据《宋史》本纪纪年）：南宋}
\noindent \eventanchor{SYD-1271-SS-144}{詔饒立削兩秩、武岡軍居住}\par

\paragraph{1272年（宋咸淳八年；公元换算据《宋史》本纪纪年）：南宋}
\noindent \eventanchor{SYD-1272-SS-145}{又詔：「有虞之世，三載考績，三考黜陟幽明}\par

\paragraph{1273年：南宋与元}
\noindent \eventanchor{SS-1273-EVENT-145}{襄阳、樊城失守}\par

\paragraph{1273年（宋咸淳九年；公元换算据《宋史》本纪纪年）：南宋}
\noindent \eventanchor{SYD-1273-SS-146}{帝覽奏，亟詔淮東制司往清口，擇利城築以備之}\par

\paragraph{1274年：南宋与元}
\noindent \eventanchor{SS-1274-EVENT-146}{元军大举南下，南宋朝廷陷入危机}\par

\paragraph{1274年（宋咸淳十年；公元换算据《宋史》本纪纪年）：南宋}
\noindent \eventanchor{SYD-1274-SS-147}{丙申，江東沙圩租米，以咸淳九年水災，詔減什四}\par

\paragraph{1275年：南宋与元}
\noindent \eventanchor{SS-1275-EVENT-147}{南宋军在丁家洲等战场失利，贾似道被罢黜}\par

\paragraph{1275年（宋德祐元年；公元换算据《宋史》本纪纪年）：南宋}
\noindent \eventanchor{SYD-1275-SS-148}{八月己亥朔，總制毛獻忠將衢州兵入衛}\par

\paragraph{1276年（宋德祐二年、元至元十三年；干支、月日待核；公元换算据双方本纪年号对照）：南宋与元}
\noindent \eventanchor{SS-1276-EVENT-148}{元军进入临安，宋恭帝降元}\par

\paragraph{1276年（宋景炎元年；公元换算据《宋史》本纪纪年）：南宋}
\noindent \eventanchor{SYD-1276-SS-149}{八月，漳州亂，命陳文龍爲閩廣宣撫使以討之}\par

\paragraph{1277年：南宋遗民与元}
\noindent \eventanchor{SS-1277-EVENT-149}{南宋遗臣、宗室和地方军民在东南继续抵抗}\par

\paragraph{1277年（宋景炎二年；公元换算据《宋史》本纪纪年）：南宋}
\noindent \eventanchor{SYD-1277-SS-150}{八月，文天祥諸將兵皆敗，乃引兵即鄒洬於永豐，洬兵亦潰}\par

\paragraph{1278年（宋祥興元年；公元换算据《宋史》本纪纪年）：南宋}
\noindent \eventanchor{SYD-1278-SS-151}{曾淵子自雷州來，以爲參知政事，廣西宣諭使}\par

\paragraph{1279年（宋祥兴二年、元至元十六年；干支、月日待核；公元换算据双方本纪年号对照）：南宋与元}
\noindent \eventanchor{SS-1279-EVENT-150}{崖山海战后南宋政权终结}\par

\paragraph{1279年（宋祥興二年；公元换算据《宋史》本纪纪年）：南宋}
\noindent \eventanchor{SYD-1279-SS-152}{大軍至中軍，會暮，且風雨，昏霧四塞，咫尺不相辨}\par
